\chapter*{Tóm tắt}
\addcontentsline{toc}{chapter}{Tóm tắt}

Báo cáo giữa kỳ này trình bày cơ sở lý thuyết cho đề tài \textbf{``Hệ thống giám sát giao thông thông minh và nhận diện biển số xe Việt Nam''}. Nội dung báo cáo tập trung vào ba phần chính:

\textbf{Chương 1: Mở đầu} -- Trình bày bối cảnh giao thông Việt Nam năm 2025, những thách thức trong việc áp dụng các hệ thống giám sát thông minh, bài toán nhận diện biển số xe, mục tiêu và phạm vi nghiên cứu của đề tài.

\textbf{Chương 2: Cơ sở lý thuyết về phát hiện và theo dõi đối tượng} -- Trình bày nền tảng toán học của bài toán phát hiện đối tượng (Object Detection), kiến trúc và cơ chế hoạt động của mô hình YOLOv8 (bao gồm Anchor-free, C2f Module, SPPF, Task Aligned Assigner), các hàm mất mát (CIoU Loss, VFL), và thuật toán theo dõi đa đối tượng ByteTrack với Bộ lọc Kalman.

\textbf{Chương 3: Cơ sở lý thuyết về nhận diện biển số xe} -- Trình bày tổng quan về hệ thống ALPR (Automatic License Plate Recognition), cấu trúc biển số xe Việt Nam, kiến trúc PaddleOCR với CRNN và CTC Loss, các kỹ thuật tiền xử lý ảnh (CLAHE, Sharpening), và cơ chế Validation \& Retry để đảm bảo kết quả OCR chính xác.

Báo cáo này đặt nền móng lý thuyết vững chắc cho việc triển khai thực nghiệm ở giai đoạn cuối kỳ, bao gồm huấn luyện mô hình, đánh giá hiệu năng và xây dựng ứng dụng hoàn chỉnh.

\vspace{0.5cm}
\textbf{Từ khóa:} YOLOv8, ByteTrack, ALPR, PaddleOCR, Kalman Filter, Object Detection, Multi-Object Tracking, License Plate Recognition.
